%% 自動生成（gen.py）
\bktocchap{3-0}{はかったデータから、きまりを\R{見}{み}つけよう}{ばねの\R{実験}{じっ|けん}から、\R{式}{しき}を\R{作}{つく}ってみよう}{7}{2}{2}{8}{J}
\bktocpart{\R{第}{だい}1\R{部}{ぶ}}{\R{関数}{かん|すう}とグラフ}{22}
\bktocchap{3-1}{\R{変}{か}わる\R{量}{りょう}と\R{変}{か}わらない\R{量}{りょう}}{ずっと\R{変}{か}わらない\R{数}{かず}と、\R{問題}{もん|だい}ごとに\R{決}{き}まる\R{数}{かず}}{7}{1}{1}{24}{C}
\bktocchap{3-2}{\R{場所}{ば|しょ}を\R{数}{かず}で\R{表}{あらわ}す：\R{座標}{ざ|ひょう}}{「どこ」を2つの\R{数}{かず}で\R{伝}{つた}えよう}{6}{1}{2}{48}{C}
\bktocchap{3-3}{\R{関数}{かん|すう}は「\R{入}{い}れると\R{出}{で}てくる」ルール}{\R{表}{ひょう}・グラフ・\R{式}{しき}は、\R{同}{おな}じことを\R{言}{い}っている}{5}{1}{2}{76}{C}
\bktocchap{3-4}{グラフをかこう・\R{読}{よ}もう}{\R{実験}{じっ|けん}の\R{点}{てん}に、なめらかな\R{線}{せん}を\R{引}{ひ}こう}{7}{1}{1}{104}{C}
\bktocchap{3-5}{グラフの\R{傾}{かたむ}きは「\R{変化}{へん|か}の\R{速}{はや}さ」}{\R{坂}{さか}が\R{急}{きゅう}なほど、\R{速}{はや}く\R{変}{か}わっている}{1}{2}{2}{128}{C}
\bktocchap{3-6}{2\R{本}{ほん}の\R{直線}{ちょく|せん}が\R{交}{まじ}わるところ：\R{連立方程式}{れん|りつ|ほう|てい|しき}}{ふたりが\R{出会}{で|あ}うのは、グラフの\R{交}{まじ}わる\R{点}{てん}}{5}{2}{2}{162}{C}
\bktocpart{\R{第}{だい}2\R{部}{ぶ}}{\R{曲}{ま}がったグラフ}{190}
\bktocchap{3-7}{\R{反比例}{はん|ぴ|れい}のグラフ}{\R{空気}{くう|き}をおしちぢめるほど、おし\R{返}{かえ}す\R{力}{ちから}は\R{強}{つよ}くなる}{1}{2}{2}{192}{C}
\bktocchap{3-8}{かっこをはずす：\R{展開}{てん|かい}と\R{平方完成}{へい|ほう|かん|せい}}{\R{式}{しき}の\R{形}{かたち}を\R{変}{か}えると、\R{最高点}{さい|こう|てん}が\R{見}{み}えてくる}{1}{2}{3}{218}{C}
\bktocchap{3-9}{\R{放物線}{ほう|ぶつ|せん}と $y=ax^2$}{\R{投}{な}げたボールは、なぜこの\R{形}{かたち}に\R{飛}{と}ぶ？}{1}{2}{3}{256}{C}
\bktocchap{3-10}{ボールはいつその\R{高}{たか}さに\R{届}{とど}く？：2\R{次方程式}{じ|ほう|てい|しき}}{\R{答}{こた}えが2つ\R{出}{で}るのは、\R{上}{のぼ}りと\R{下}{くだ}りがあるから}{1}{2}{3}{290}{C}
\bktocpart{\R{第}{だい}3\R{部}{ぶ}}{\R{倍々}{ばい|ばい}とけた}{326}
\bktocchap{3-11}{\R{指数関数}{し|すう|かん|すう}：いまある\R{量}{りょう}に\R{比例}{ひ|れい}してふえる・へる}{\R{倍々}{ばい|ばい}でふえると、あっという\R{間}{ま}に\R{大}{おお}きくなる}{3}{2}{3}{328}{C}
\bktocchap{3-12}{\R{対数}{たい|すう}：けたを\R{数}{かぞ}える\R{道具}{どう|ぐ}}{\R{音}{おと}の\R{大}{おお}きさの「デシベル」は、けたの\R{数}{かぞ}え\R{方}{かた}}{3}{2}{3}{358}{C}
\bktocpart{\R{第}{だい}4\R{部}{ぶ}}{データと\R{統計}{とう|けい}}{386}
\bktocchap{3-13}{データのちらばり：ヒストグラムと\R{平均}{へい|きん}}{\R{何回}{なん|かい}もはかって、\R{本当}{ほん|とう}の\R{値}{あたい}に\R{近}{ちか}づこう}{7}{2}{2}{388}{C}
\bktocchap{3-14}{ばらつきを1つの\R{数}{かず}で：\R{標準偏差}{ひょう|じゅん|へん|さ}}{「だいたいこのくらいずれる」を\R{数}{かず}で\R{言}{い}う}{7}{3}{3}{416}{C}
\bktocchap{3-15}{2つの\R{量}{りょう}の\R{関係}{かん|けい}：\R{散布図}{さん|ぷ|ず}と\R{相関}{そう|かん}}{\R{背}{せ}の\R{高}{たか}い\R{人}{ひと}は、\R{足}{あし}も\R{速}{はや}い？}{7}{3}{3}{442}{C}
\bktocchap{3-16}{データに\R{直線}{ちょく|せん}を\R{引}{ひ}く：\R{最小二乗法}{さい|しょう|に|じょう|ほう}}{ずれがいちばん\R{小}{ちい}さくなる\R{線}{せん}は、どれ？}{7}{3}{3}{474}{C}
\bktocchap{3-17}{\R{偶然}{ぐう|ぜん}のばらつき：\R{確率}{かく|りつ}と\R{正規分布}{せい|き|ぶん|ぷ}}{サイコロをたくさん\R{足}{た}すと、つりがねの\R{形}{かたち}に}{7}{3}{3}{498}{C}
\par\vspace{2mm}
\bktocbacklist{\bkbi{こたえ}{528}\bkbi{\R{記号}{き|ごう}カード}{544}\bkbi{よく\R{使}{つか}う\R{定数}{てい|すう}・\R{単位}{たん|い}・\R{公式}{こう|しき}}{548}\bkbi{できるようになったことチェック\R{表}{ひょう}}{555}\bkbi{つぎの\R{巻}{かん}へ}{558}\bkbi{さくいん}{561}}
